% ROZDZIAŁ 21 — RIESZ–MARKOW–KAKUTANI
\chapter{Twierdzenie Riesza--Markowa--Kakutaniego}
\label{ch:riesz-markov-21}

\providecommand{\atlasref}[2]{#2}


\section{Sformułowanie i rola twierdzenia}
Niech $X$ będzie lokalnie zwartą przestrzenią Hausdorffa,
a $\mathbb K\in\{\mathbb R,\mathbb C\}$.
Przestrzeń $C_c(X;\mathbb K)$ i pojęcie zwartego nośnika
wprowadzono w rozdziale~\ref{ch:radon-measures-20}.
Przez $C_0(X;\mathbb K)$ oznaczamy przestrzeń funkcji ciągłych
\emph{znikających w nieskończoności}, to znaczy takich, że
dla każdego $\varepsilon>0$ zbiór
$\{x:|f(x)|\ge\varepsilon\}$ jest zwarty.
Przestrzeń $C_0(X;\mathbb K)$ z normą supremum
jest przestrzenią Banacha. Najpierw zauważmy, że
$C_c(X;\mathbb K)\subseteq C_0(X;\mathbb K)$,
gdyż dodatnie zbiory poziomicowe modułu funkcji
o zwartym nośniku są domkniętymi podzbiorami tego
nośnika. Z nierówności trójkąta
$\{|f+g|\ge\varepsilon\}\subseteq
\{|f|\ge\varepsilon/2\}\cup\{|g|\ge\varepsilon/2\}$
wynika zwartość zbioru po lewej, gdyż jest domknięty.
Dla skalara $a\ne0$ mamy
$\{|af|\ge\varepsilon\}=\{|f|\ge\varepsilon/|a|\}$,
a dla $a=0$ zbiór jest pusty. Otrzymujemy przestrzeń
wektorową. Każda jej funkcja jest ograniczona:
jeśli zwarty zbiór $\{|f|\ge1\}$ jest niepusty,
moduł osiąga na nim skończone maksimum, a poza nim
jest mniejszy od $1$; jeśli jest pusty, moduł jest
wszędzie mniejszy od $1$. Dla pustego $X$
przyjmujemy normę jedynej funkcji równą zero. Dla ciągu Cauchy'ego $(f_n)$ w normie supremum
każdy ciąg $(f_n(x))$ jest ciągiem Cauchy'ego w
zupełnym ciele $\mathbb K$, ma więc granicę $f(x)$.
Zbieżność do $f$ jest jednostajna, ponieważ po ustaleniu $\varepsilon>0$ w nierówności $\|f_n-f_m\|_\infty\le\varepsilon$ dla $n,m$ dostatecznie dużych przechodzimy z $m\to\infty$. Granica $f$ jest ciągła. Jeżeli $\|f-f_n\|_\infty<\varepsilon/2$, to domknięty zbiór $\{|f|\ge\varepsilon\}$ zawiera się w zwartym zbiorze $\{|f_n|\ge\varepsilon/2\}$, jest więc zwarty. To daje zupełność.
Zupełność $C(K;\mathbb K)$ dla zwartej przestrzeni Hausdorffa $K$ omówiono w
\atlasref{FA-12::prop:ck-banach12}{FA-12: zupełność przestrzeni funkcji ciągłych}. Dla dowolnej lokalnie zwartej przestrzeni Hausdorffa udowodnimy gęstość
$C_c(X;\mathbb K)$ w $C_0(X;\mathbb K)$.

Konstrukcję miary z dodatniego funkcjonału całkowego
na kracie funkcji przedstawiono z pełnym dowodem w
\atlasref{ME-32::it:thm:daniell-stone}{ME-32: twierdzenie Daniella--Stone'a}
oraz w wersji bez jedynki w
\atlasref{ME-32::it:thm:daniell-stone-general}{ME-32: warunek Stone'a i ogólna reprezentacja}.
Ciągłość porządkowa oznacza tam malejącą zbieżność
punktową funkcji do zera, a nie tylko zerowe infimum
w kracie. W przypadku $C_c(X;\mathbb R)$ ciągłość ta
wynika z twierdzenia Diniego na zwartym nośniku
pierwszego wyrazu i oszacowania przez funkcję odcinającą.
Dla lokalnie zwartej ośrodkowej przestrzeni metrycznej
przeliczalne pokrycie przez zbiory niezerowości funkcji
z $C_c(X)$ wynika z $\sigma$-zwartości.
W \atlasref{ME-32::it:ex:daniell-cc}{ME-32: zastosowanie do funkcji o zwartym nośniku}
wyprowadzono w ten sposób reprezentację w przypadku metrycznym.
Poniżej podajemy konstrukcję ogólną, a następnie alternatywny
dowód metryczny. Porównanie dziedzin i wartości miar przedstawia
\atlasref{ME-32::it:thm:daniell-radon-comparison}{ME-32: porównanie miar Daniella i Radona}.

\begin{twierdzenie}[Reprezentacja dodatnich funkcjonałów]
\label{thm:riesz-positive-21}
Niech $X$ będzie \emph{lokalnie zwartą przestrzenią Hausdorffa}, a $L:C_c(X;\mathbb R)\to\mathbb R$ funkcjonałem
liniowym takim, że $L(f)\ge0$ dla $f\ge0$.
Istnieje dokładnie jedna miara Radona $\mu$ na $X$ spełniająca
\begin{equation}\label{eq:riesz-representation-21}
 L(f)=\int_X f(x)\,\mathrm d\mu(x),\qquad f\in C_c(X;\mathbb R).
\end{equation}
Przy czym $L$ nie musi być ograniczony w normie supremum na
całej przestrzeni $C_c(X)$.
\end{twierdzenie}

\begin{uwaga}[Zakres twierdzenia]
Wersja ogólna wykorzystuje konwencję Radona z definicji \ref{def:radon-measure-20}, czyli regularność wewnętrzną na otwartych. Nie wymagamy $\sigma$-skończoności miary ani metryczności przestrzeni. Dodatkowa regularność na wszystkich zbiorach borelowskich wynika z lematu \ref{lem:radon-sigmafinite-inner-20}, gdy $X$ jest $\sigma$-zwarta. W szczególności dotyczy to lokalnie zwartych ośrodkowych przestrzeni metrycznych. Poniżej podajemy także wcześniejszą alternatywną konstrukcję metryczną.
\end{uwaga}

\section{Dowód ogólny: lokalnie zwarte przestrzenie Hausdorffa}
\begin{lemat}[Lokalne ograniczenie funkcjonału]\label{it:lem:radon-local-bound}
Niech $X$ będzie lokalnie zwartą przestrzenią Hausdorffa
i niech $L:C_c(X;\mathbb R)\to\mathbb R$ będzie
dodatnim funkcjonałem liniowym.
Dla każdego zwartego $H\subseteq X$ istnieje $\eta\in C_c(X)$, $0\le\eta\le1$, $\eta|_H=1$, taka, że
\[
 |L(f)|\le L(\eta)\|f\|_\infty
 \quad\text{dla }f\in C_c(X),\quad\operatorname{supp}f\subset H.
\]
\end{lemat}
\begin{proof}
Lemat \ref{lem:cutoff-radon-20} zastosowany do $H\subset X$ daje $\eta$. Punktowo $-\|f\|_\infty\eta\le f\le\|f\|_\infty\eta$. Dodatniość i liniowość $L$ dają oszacowanie. Nie używamy tutaj istnienia miary.
\end{proof}

\begin{proof}[Dowód twierdzenia \ref{thm:riesz-positive-21} w ogólności]
Jeśli $X=\varnothing$, jedyną miarą jest miara zerowa i teza zachodzi. W dalszej części nie zakładamy, że $X$ ma przeliczalne pokrycie zbiorami zwartymi.

\emph{Krok 1: zawartość zbiorów otwartych.}
Dla otwartego $U\subset X$ definiujemy
\begin{equation}\label{it:eq:radon-open-content-general}
 a(U)=\sup\{L(f): f\in C_c(X),\ 0\le f\le1,\
 \operatorname{supp}f\subset U\}.
\end{equation}
Funkcja zerowa jest zawsze dopuszczalna, więc $a$ przyjmuje wartości w $[0,\infty]$, jest monotoniczna oraz $a(\varnothing)=0$. Dodatniość i liniowość $L$ implikują monotoniczność na $C_c(X)$, ponieważ dla $f\le g$ mamy $L(g)-L(f)=L(g-f)\ge0$. Jeśli $\overline U$ jest zwarty, z lematu \ref{it:lem:radon-local-bound} mamy $a(U)\le L(\eta)<\infty$. Dla rozłącznych otwartych $U,V$ i funkcji dopuszczalnych $f,g$ odpowiednio dla $U,V$ mamy $0\le f+g\le1$ i $\operatorname{supp}(f+g)\subset U\cup V$. Zatem $a(U\cup V)\ge L(f)+L(g)$. Biorąc niezależne suprema, otrzymujemy $a(U\cup V)\ge a(U)+a(V)$, również jeśli jeden ze składników jest nieskończony.

\emph{Krok 2: podaddytywność na otwartych.}
Jeśli $U\subset\bigcup_{j\ge1}U_j$ dla otwartych $U_j$, ustalmy funkcję dopuszczalną $f$ dla $U$. Jeśli $f=0$, oszacowanie jest natychmiastowe. W przeciwnym przypadku dla każdego $x$ ze zwartego $H=\operatorname{supp}f$ wybieramy $j$ i otwarte otoczenie $V_x$ o zwartym domknięciu zawartym w $U_j$. Istnienie wynika z regularności i lokalnej zwartości, jak w dowodzie lematu \ref{lem:cutoff-radon-20}. Ze zwartości $H$ otrzymujemy skończone pokrycie $V_1,\ldots,V_m$ z $\overline V_i\subset U_{j_i}$. Lemat odcinający daje $\psi_i\in C_c(X)$, $0\le\psi_i\le1$, $\psi_i=1$ na $\overline V_i$ i $\operatorname{supp}\psi_i\subset U_{j_i}$. Wówczas $\psi=\sum_i\psi_i$ jest dodatnia na $H$. Definiujemy
\[
 f_i(x)=
 \begin{cases}
 f(x)\psi_i(x)/\psi(x),&\psi(x)>0,\\
 0,&\psi(x)=0.
 \end{cases}
\]
Na otwartym zbiorze $\{\psi>0\}$ funkcja $f_i$ jest ciągła. W każdym punkcie $x$ z $\psi(x)=0$ mamy $x\notin H$, a zatem $f=0$ w pewnym otoczeniu $x$, co daje ciągłość także tam. Ponadto $0\le f_i\le f$ i zbiór punktów, w których $f_i\ne0$, zawiera się w przecięciu zbiorów, w których $f\ne0$ i $\psi_i\ne0$. Po wzięciu domknięć otrzymujemy $\operatorname{supp}f_i\subseteq H\cap\operatorname{supp}\psi_i\subseteq U_{j_i}$. Ponieważ $H$ jest zwarty, każdy $f_i$ ma zwarty nośnik. Na $\{\psi>0\}$ suma ilorazów wynosi $f$, a na $\{\psi=0\}$ mamy $f=0$; stąd $\sum_i f_i=f$ na całym $X$. Grupujemy powtarzające się indeksy: dla każdego występującego $j$ kładziemy $g_j=\sum_{i:j_i=j}f_i$. Wówczas $0\le g_j\le f\le1$ i $\operatorname{supp}g_j\subset U_j$. Z liniowości $L(f)=\sum_jL(g_j)\le\sum_{j\ge1}a(U_j)$, gdzie lewa suma jest skończona. Biorąc supremum po $f$, otrzymujemy $a(U)\le\sum_ja(U_j)$. W połączeniu z krokiem 1 daje to addytywność na dwóch rozłącznych zbiorach otwartych.

\emph{Krok 3: miara zewnętrzna i zwarta aproksymacja przed mierzalnością.}
Określamy dla dowolnego $E\subset X$
\begin{equation}\label{it:eq:radon-outer-general}
 m^*(E)=\inf\{a(U):E\subset U,\ U\text{ otwarty}\}.
\end{equation}
Mamy $m^*(\varnothing)=0$ oraz monotoniczność. Jeśli $E\subset\bigcup_jE_j$ i $\sum_jm^*(E_j)<\infty$, dla $\varepsilon>0$ wybieramy otwarte $U_j\supset E_j$ z $a(U_j)<m^*(E_j)+\varepsilon2^{-j}$. Wtedy
$m^*(E)\le a(\bigcup_jU_j)\le\sum_jm^*(E_j)+\varepsilon$.
Przechodząc z $\varepsilon\downarrow0$, otrzymujemy przeliczalną podaddytywność. Dla nieskończonej sumy podaddytywność jest automatyczna. Stąd $m^*$ jest miarą zewnętrzną. Dla otwartego $U$ zachodzi $m^*(U)=a(U)$: sam zbiór $U$ jest dopuszczalny w infimum, a każdy jego otwarty nadzbiór $V$ spełnia $a(V)\ge a(U)$. Każdy zwarty $H$ ma otwarte otoczenie o zwartym domknięciu (skończona suma takich otoczeń punktów $H$), więc $m^*(H)<\infty$.

Jeżeli $f$ jest dopuszczalna dla $U$, to dla $H=\operatorname{supp}f\subset U$ i każdego otwartego $V\supset H$ zachodzi $a(V)\ge L(f)$. Stąd $m^*(H)\ge L(f)$. Biorąc supremum po $f$, dostajemy
\begin{equation}\label{it:eq:radon-outer-inner-general}
 m^*(U)=a(U)=\sup\{m^*(H):H\subset U,\ H\text{ zwarty}\}.
\end{equation}
Odwrotna nierówność wynika z monotoniczności $m^*$; argument działa także dla nieskończonego $a(U)$. Na tym etapie nie zakładamy mierzalności $H$ ani $U$.

\emph{Krok 4: kryterium Carathéodory'ego dla otwartych.}
Ustalmy otwarty $U\subset X$ oraz dowolny $E\subset X$. Jeśli $m^*(E)=+\infty$, podaddytywność wymusza, że suma $m^*(E\cap U)+m^*(E\setminus U)$ jest nieskończona, więc kryterium zachodzi. Niech $m^*(E)<\infty$ i $\varepsilon>0$. Wybieramy otwarty $W\supset E$ z $a(W)<m^*(E)+\varepsilon$. Mamy $a(W\cap U)\le a(W)<\infty$. Z \eqref{it:eq:radon-outer-inner-general} wybieramy zwarty $H\subset W\cap U$ z
$m^*(H)>a(W\cap U)-\varepsilon$ (jeśli $a(W\cap U)=0$, wystarczy $H=\varnothing$). Z lematu \atlasref{GT-14::lem:topog-baire-zwarte-domkniecie}{GT-14: otoczenie o zwartym domknięciu} i skończonego podpokrycia zwartego $H$ wybieramy otwarty $V$ taki, że $H\subseteq V\subseteq\overline V\subseteq W\cap U$ i $\overline V$ jest zwarty; dla pustego $H$ można przyjąć $V=\varnothing$. Wówczas $V$ i $W\setminus\overline V$ są rozłącznymi otwartymi podzbiorami $W$, a $E\setminus U\subset W\setminus\overline V$. Z kroku 1 i monotoniczności mamy
\begin{align*}
 m^*(E)+\varepsilon
 &>a(W)\ge a(V)+a(W\setminus\overline V)\\
 &\ge m^*(H)+m^*(E\setminus U)\\
 &>a(W\cap U)-\varepsilon+m^*(E\setminus U)\\
 &\ge m^*(E\cap U)+m^*(E\setminus U)-\varepsilon .
\end{align*}
Zatem $m^*(E\cap U)+m^*(E\setminus U)<m^*(E)+2\varepsilon$. Z dowolności $\varepsilon$ otrzymujemy nierówność potrzebną w kryterium Carathéodory'ego; przeciwna jest podaddytywnością. Wszystkie zbiory otwarte są więc mierzalne względem $m^*$. Twierdzenie \ref{thm:caratheodory} daje mierzalność wszystkich zbiorów borelowskich i przeliczalną addytywność $\mu=m^*|_{\mathcal B(X)}$. Z kroków 3--4 wynika skończoność na zwartych, regularność zewnętrzna i regularność wewnętrzna na otwartych. Zatem $\mu$ jest miarą Radona według definicji \ref{def:radon-measure-20}.

\emph{Krok 5: identyfikacja funkcjonału z całką.}
Niech $f\in C_c(X)$, $f\ge0$, $f\ne0$. Oznaczmy $H=\operatorname{supp}f$, $M=\|f\|_\infty>0$, wybierzmy $\eta$ z lematu \ref{it:lem:radon-local-bound}, a dla $N\ge1$ połóżmy $\delta=M/N$ i $U_i=\{f>i\delta\}$, $0\le i\le N$. Wszystkie $U_i$ są otwarte, zawarte w $H$, zatem mają skończoną miarę. Dla $0\le i<N$ zbiór $H_i=\{f\ge(i+1)\delta\}$ jest zwartym podzbiorem $U_i$. Z lematu \ref{lem:cutoff-radon-20} wybieramy $h_i\in C_c(X)$, $0\le h_i\le1$, $h_i|_{H_i}=1$ i $\operatorname{supp}h_i\subset U_i$. Punktowo $f\le\delta\eta+\delta\sum_{i=0}^{N-1}h_i$: dla $x\in H$ i $q=\lfloor f(x)/\delta\rfloor$ mamy $0\le q\le N$, $\eta(x)=1$ oraz $h_i(x)=1$ dla $0\le i<q$. Zatem prawa strona jest nie mniejsza od $(q+1)\delta\ge f(x)$. Poza $H$ lewa strona jest zerowa, a prawa nieujemna. Dodatniość daje
\[
 L(f)\le\delta L(\eta)+\delta\sum_{i=0}^{N-1}\mu(U_i).
\]
Dla dowolnych dopuszczalnych funkcji $g_i$ dla $U_i$, $1\le i\le N$, mamy $\delta\sum_i g_i\le f$, ponieważ liczba progów $i\delta<f(x)$ nie przekracza $f(x)/\delta$. Z dodatniości i wzięcia niezależnych supremów (skończona liczba skończonych supremów) wynika $\delta\sum_{i=1}^N\mu(U_i)\le L(f)$. Dla $0\le y\le M=N\delta$ mamy
$\delta\sum_{i=1}^{N}\chi_{(i\delta,\infty)}(y)\le y
\le\delta\sum_{i=0}^{N-1}\chi_{(i\delta,\infty)}(y)$:
przy $y=0$ obie sumy są zerowe; dla
$j\delta<y\le(j+1)\delta$, $0\le j<N$,
dolna suma wynosi $j\delta$, a górna $(j+1)\delta$.
Stosując te nierówności do $y=f(x)$ i całkując
mierzalne funkcje proste, otrzymujemy
\[
 \delta\sum_{i=1}^{N}\mu(U_i)
 \le\int_X f\,\mathrm d\mu
 \le\delta\sum_{i=0}^{N-1}\mu(U_i).
\]
Różnica sum wynosi $\delta\mu(U_0)\le\delta\mu(H)$, ponieważ $U_N=\varnothing$. Stąd
\[
 \left|L(f)-\int_X f\,\mathrm d\mu\right|
 \le\delta\bigl(L(\eta)+\mu(H)\bigr)\longrightarrow0.
\]
Zatem zachodzi reprezentacja dla $f\ge0$ (dla $f=0$ obie strony są zerowe). Dla dowolnej rzeczywistej $f\in C_c(X)$ stosujemy wynik do $f^+$ i $f^-$, które również mają zwarty nośnik, i odejmujemy skończone całki. Jednoznaczność wynika z twierdzenia \ref{thm:uniqueness-cc-20}. 
\end{proof}



\begin{wniosek}[Odzyskiwanie miary zbioru zwartego z funkcjonału]
\label{it:cor:radon-compact-inf}
Niech $L$ i $\mu$ spełniają tezę twierdzenia
\ref{thm:riesz-positive-21}. Dla każdego zwartego
$K\subseteq X$ zachodzi
\[
 \mu(K)=\inf\{L(f):f\in C_c(X),\
             0\le f\le1,\ f|_K=1\}.
\]
Wraz ze wzorem \eqref{eq:cc-open-sup-20} daje to
dwie formuły odzyskiwania miary bez znajomości
jej wartości na dowolnych zbiorach borelowskich.
\end{wniosek}
\begin{proof}
Każda dopuszczalna funkcja $f$ spełnia
$\chi_K\le f$. Z reprezentacji i monotoniczności
całki otrzymujemy $\mu(K)\le L(f)$.
Dla $\varepsilon>0$ regularność zewnętrzna
i skończoność $\mu(K)$ dają otwarty $U\supseteq K$
z $\mu(U)<\mu(K)+\varepsilon$.
Lemat \ref{lem:cutoff-radon-20} daje dopuszczalną
funkcję $f$ o nośniku zawartym w $U$.
Mamy $f\le\chi_U$, zatem
$L(f)=\int_X f\,\mathrm d\mu\le\mu(U)
<\mu(K)+\varepsilon$.
Dowolność $\varepsilon$ dowodzi równości.
Dla $K=\varnothing$ można też bezpośrednio
przyjąć $f=0$.
\end{proof}

\section{Lemat o metrycznych miarach zewnętrznych}
\begin{definicja}[Miara zewnętrzna metryczna]
\label{def:metric-outer-21}
Niech $(Z,d)$ będzie przestrzenią metryczną.
Miara zewnętrzna $m^*:\mathcal P(Z)\to[0,\infty]$
jest \emph{metryczna}, gdy dla dowolnych $A,B\subseteq Z$
o dodatniej odległości $d(A,B)>0$ zachodzi
$m^*(A\cup B)=m^*(A)+m^*(B)$.
Dla niepustych zbiorów przyjmujemy
$d(A,B)=\inf\{d(x,y):x\in A,\ y\in B\}$,
a jeśli co najmniej jeden ze zbiorów jest pusty,
przyjmujemy $d(A,B)=+\infty$.
\end{definicja}

\begin{lemat}[Borelowska mierzalność]
\label{lem:metric-outer-borel-21}
Jeżeli $m^*$ jest metryczną miarą zewnętrzną na $(Z,d)$,
to każdy zbiór borelowski w $Z$ jest mierzalny w sensie
Carathéodory'ego.
\end{lemat}
\begin{proof}
Wystarczy wykazać mierzalność zbiorów domkniętych, gdyż
rodzina zbiorów mierzalnych Carathéodory'ego jest
$\sigma$-algebrą (twierdzenie \atlasref{ME-03::thm:caratheodory}{ME-03: twierdzenie Carathéodory'ego}).
Niech $F\subseteq Z$ będzie domknięty i $E\subseteq Z$.
Musimy dowieść
$m^*(E)\ge m^*(E\cap F)+m^*(E\setminus F)$;
nierówność odwrotna wynika z podaddytywności.
Jeśli $F=\varnothing$ lub $F=Z$, mierzalność jest oczywista.
Jeśli $m^*(E)=+\infty$, podaddytywność wymusza nieskończoną sumę po prawej, więc kryterium zachodzi. Załóżmy
$m^*(E)<\infty$ i oznaczmy $r(x)=d(x,F)$.
Dla każdego $n\ge1$ zbiory $E\cap F$ oraz
$E\cap\{r>1/n\}$ mają dodatnią odległość, gdyż
$r$ jest $1$-lipschitzowska. Zatem
\begin{equation}\label{eq:metric-outer-sep-21}
 m^*(E)\ge m^*(E\cap F)+m^*(E\cap\{r>1/n\}).
\end{equation}
Wykażemy, że drugi składnik po prawej zbiega do
$m^*(E\setminus F)$; nie wolno tutaj korzystać z ciągłości
z dołu, ponieważ $m^*$ jest dopiero miarą zewnętrzną.
Niech $A_0=E\cap\{r>1\}$ oraz
$A_j=E\cap\{1/(j+1)<r\le1/j\}$, $j\ge1$.
Dla $j\ge1$ i $k\ge j+2$, jeśli $x\in A_j$,
$y\in A_k$, to z $1$-lipschitzowskości $r$ wynika
\[
 d(x,y)\ge r(x)-r(y)>
 \frac1{j+1}-\frac1{j+2}>0.
\]
Zatem $d(A_j,A_k)\ge1/(j+1)-1/(j+2)>0$.
Także $d(A_0,A_k)\ge1/2$ dla $k\ge2$.
Dla skończonej rodziny warstw tej samej parzystości
odległość każdej warstwy od sumy pozostałych
jest nie mniejsza od minimum skończenie wielu
dodatnich odległości (puste warstwy można pominąć).
Metryczność i indukcja dają więc
$\sum m^*(A_j)=m^*(\bigcup A_j)\le m^*(E)$.
W konsekwencji oba szeregi, po indeksach parzystych
i nieparzystych, są zbieżne; zatem
$\sum_{j\ge1}m^*(A_j)<\infty$. Ponieważ
\[
 E\cap\{0<r\le1/n\}\subseteq\bigcup_{j\ge n}A_j,
\]
miara zewnętrzna części po lewej dąży do zera.
Z podaddytywności
\[
 m^*(E\setminus F)\le
 m^*(E\cap\{r>1/n\})+
 m^*(E\cap\{0<r\le1/n\}),
\]
a z monotoniczności wynika przeciwna nierówność
$m^*(E\cap\{r>1/n\})\le m^*(E\setminus F)$.
Przechodząc do granicy w \eqref{eq:metric-outer-sep-21},
otrzymujemy żądane kryterium Carathéodory'ego.
\end{proof}

\section{Przypadek zwarty: konstrukcja od podstaw}

\noindent\emph{Powiązania tematyczne w Bibliotece matematycznej:} \atlasref{FA-12::thm:riesz-markov12}{FA-12: reprezentacja funkcjonałów na przestrzeni funkcji ciągłych}.

\begin{twierdzenie}[Reprezentacja na zwartej przestrzeni metrycznej]
\label{thm:compact-riesz-21}
Niech $K$ będzie niepustą zwartą przestrzenią metryczną,
a $L:C(K;\mathbb R)\to\mathbb R$ dodatnim funkcjonałem
liniowym. Istnieje dokładnie jedna skończona regularna
miara borelowska $\mu$ na $K$ taka, że
$L(f)=\int_Kf\,\mathrm d\mu$ dla wszystkich $f\in C(K)$.
Ponadto $\mu(K)=L(\chi_K)$ i
$\|L\|=L(\chi_K)$.
\end{twierdzenie}
\begin{proof}
\emph{Krok 1: lokalne ograniczenie funkcjonału.}
Z $-\|f\|_\infty\chi_K\le f\le\|f\|_\infty\chi_K$
wynika $|L(f)|\le L(\chi_K)\|f\|_\infty$.
Dla $f=\chi_K$ mamy równość w normie.

\emph{Krok 2: przypisanie liczb zbiorom otwartym.}
Dla zbioru otwartego $U\subseteq K$ definiujemy
\begin{equation}\label{eq:open-content-21}
 a(U)=\sup\{L(f): f\in C(K),\ 0\le f\le1,\
          \operatorname{supp}f\subseteq U\}.
\end{equation}
W szczególności $a(\varnothing)=0$,
$0\le a(U)\le L(\chi_K)$, $a(K)=L(\chi_K)$,
a $a$ jest monotoniczna. Jeżeli $U,V$ są rozłączne
i otwarte, to $a(U\cup V)\ge a(U)+a(V)$:
wybieramy funkcje $f,g$ dopuszczalne odpowiednio dla
$U,V$; ponieważ ich nośniki są rozłączne,
$0\le f+g\le1$ i $L(f+g)=L(f)+L(g)$.

\emph{Krok 3: przeliczalna podaddytywność na otwartych.}
Jeżeli $U\subseteq\bigcup_{j\ge1}U_j$ i $f$ jest
dopuszczalna w \eqref{eq:open-content-21}, przypadek
$f=0$ daje od razu $L(f)=0\le\sum_j a(U_j)$.
W przeciwnym razie $\operatorname{supp}f$ ma skończone pokrycie
$U_{j_1},\ldots,U_{j_N}$. Zbudujmy ciągłe funkcje
$\theta_i\ge0$ o nośnikach zawartych w $U_{j_i}$,
takie, że $\sum_i\theta_i=1$ na
$H:=\operatorname{supp}f$. Dla każdego $x\in H$
wybieramy otwarty zbiór $V_x$ zawierający $x$,
którego domknięcie w $K$ mieści się w pewnym $U_{j_i}$.
Ze zwartości $H$ wybieramy skończone pokrycie
takimi $V_x$; po ewentualnym powtórzeniu indeksów
oznaczmy je przez $V_i$, tak aby
$\overline V_i\subseteq U_{j_i}$.
Lemat~\ref{lem:cutoff-radon-20} daje funkcje
$\psi_i\in C(K)$ spełniające
$0\le\psi_i\le1$, $\psi_i=1$ na $\overline V_i$ i
$\operatorname{supp}\psi_i\subseteq U_{j_i}$.
Dla $\psi=\sum_i\psi_i$ mamy $\psi>0$ na $H$.
Ponownie stosujemy lemat o funkcji odcinającej,
uzyskując $\chi\in C(K)$ taką, że $0\le\chi\le1$,
$\chi=1$ na $H$ oraz
$\operatorname{supp}\chi\subseteq\{\psi>0\}$.
Wówczas $\theta_i=\chi\psi_i/\psi$ na $\{\psi>0\}$
i $\theta_i=0$ poza tym zbiorem. Funkcje
$\theta_i$ są ciągłe również na brzegu $\{\psi>0\}$,
ponieważ nośnik $\chi$ jest w nim zawarty,
a ich suma jest równa $\chi=1$ na $H$. Każde $f\theta_i$ jest nieujemne,
nie przekracza $1$ i ma nośnik w $U_{j_i}$.
Jeżeli indeks $j_i$ występuje więcej niż raz, grupujemy odpowiadające mu funkcje: $h_j=\sum_{i:j_i=j}f\theta_i$. Mamy $0\le h_j\le f\le1$ (ponieważ suma wszystkich składników jest równa $f$), a jego nośnik zawiera się w $U_j$, jako skończona suma odpowiednich nośników. Wobec tego
$L(f)=\sum_j L(h_j)\le\sum_{j\ge1}a(U_j)$.
Po wzięciu supremum otrzymujemy
$a(U)\le\sum_j a(U_j)$.

\emph{Krok 4: miara zewnętrzna.}
Dla dowolnego $A\subseteq K$ definiujemy
\begin{equation}\label{eq:riesz-outer-21}
 m^*(A)=\inf\{a(U):A\subseteq U,\ U\text{ otwarty}\}.
\end{equation}
Jeżeli $A\subseteq\bigcup_j A_j$, dobieramy otwarte
$U_j\supseteq A_j$ z
$a(U_j)\le m^*(A_j)+\varepsilon2^{-j}$ (przyjmujemy
osobno trywialny przypadek nieskończonej sumy).
Krok 3 daje
$m^*(A)\le a(\bigcup_jU_j)\le\sum_jm^*(A_j)+\varepsilon$.
Przechodząc z $\varepsilon\downarrow0$ otrzymujemy
przeliczalną podaddytywność. Ponadto
$m^*(\varnothing)=0$ z $a(\varnothing)=0$,
a monotoniczność wynika z zawierania się rodzin
dopuszczalnych otwartych nadzbiorów. Zatem $m^*$
jest miarą zewnętrzną.
Dla zbioru otwartego $U$ mamy $m^*(U)=a(U)$, ponieważ
$U\subseteq V$ implikuje $a(U)\le a(V)$.
Jeśli $A$ lub $B$ jest pusty, metryczna addytywność zachodzi wprost. W przeciwnym razie, jeśli $d(A,B)=\delta>0$, to dla każdego otwartego
$W\supseteq A\cup B$ zbiory
$W_A=W\cap\{x:d(x,A)<\delta/3\}$ i
$W_B=W\cap\{x:d(x,B)<\delta/3\}$ są otwarte,
rozłączne i zawierają odpowiednio $A,B$. Stąd
$a(W)\ge a(W_A)+a(W_B)\ge m^*(A)+m^*(B)$.
Rozłączność $W_A,W_B$ wynika z nierówności
$d(A,B)\le d(x,A)+d(x,B)$, która zachodzi dla
każdego $x$ na mocy nierówności trójkąta i wzięcia
infimum po punktach obu zbiorów. Gdyby $x$ należał
do obu otoczeń, otrzymalibyśmy $\delta<2\delta/3$.
Biorąc infimum po $W$ otrzymujemy
$m^*(A\cup B)\ge m^*(A)+m^*(B)$; nierówność
przeciwna wynika z podaddytywności. To dowodzi
metryczności $m^*$.
Lemat~\ref{lem:metric-outer-borel-21} zapewnia, że
wszystkie zbiory borelowskie są mierzalne. Ograniczenie
$\mu=m^*|_{\mathcal B(K)}$ jest więc miarą, a
$\mu(K)=L(\chi_K)<\infty$.

\emph{Krok 5: regularność.}
Regularność zewnętrzna wynika z \eqref{eq:riesz-outer-21}.
Dla otwartego $U$ i $\varepsilon>0$ dobieramy
$0\le f\le1$ z $\operatorname{supp}f\subseteq U$ i
$L(f)>\mu(U)-\varepsilon$. Przyjmijmy
$H=\operatorname{supp}f$, który jest zwarty.
Każdy otwarty $V\supseteq H$ spełnia
$L(f)\le a(V)$, zatem
$\mu(H)=\inf_{V\supseteq H}a(V)\ge L(f)$.
To dowodzi regularności wewnętrznej dla otwartych.
Dla dowolnego borelowskiego $E$ i $\varepsilon>0$
wybieramy otwarty $U\supseteq E$ z
$\mu(U\setminus E)<\varepsilon$.
Następnie wybieramy zwarty $H\subseteq U$ z
$\mu(U\setminus H)<\varepsilon$ oraz otwarty
$W\supseteq U\setminus E$ z
$\mu(W)<\mu(U\setminus E)+\varepsilon$.
Zbiór $H\setminus W$ jest zwarty, zawiera się w $E$,
a $\mu(E\setminus(H\setminus W))<3\varepsilon$,
ponieważ $E\setminus(H\setminus W)\subseteq(U\setminus H)\cup W$
i $\mu(W)<2\varepsilon$.
Otrzymujemy regularność wewnętrzną względem zwartych.

\emph{Krok 6: identyfikacja całki z $L$.}
Niech $f\in C(K)$, $f\ge0$. Dla $f=0$ równość jest natychmiastowa; w pozostałym przypadku $M=\|f\|_\infty>0$.
Dla $N\ge1$ połóżmy $\delta=M/N$ i
$U_i=\{f>i\delta\}$, $i=0,\ldots,N$.
Dla $i=0,\ldots,N-1$ zbiór
$H_i=\{f\ge(i+1)\delta\}$ jest zwarty i zawiera się
w $U_i$. Istnieje $h_i\in C(K)$, $0\le h_i\le1$,
$h_i=1$ na $H_i$, $\operatorname{supp}h_i\subseteq U_i$;
wynika to z lematu o funkcji odcinającej
\ref{lem:cutoff-radon-20}. Punktowo
$f\le\delta\chi_K+\delta\sum_{i=0}^{N-1}h_i$.
Z dodatniości $L$ dostajemy
\begin{equation}\label{eq:riesz-upper-sum-21}
 L(f)\le\delta L(\chi_K)+
           \delta\sum_{i=0}^{N-1}\mu(U_i).
\end{equation}
Z drugiej strony dla dowolnych funkcji
$g_i\in C(K)$, $0\le g_i\le1$,
$\operatorname{supp}g_i\subseteq U_i$ ($1\le i\le N$),
zachodzi $\delta\sum_{i=1}^{N}g_i\le f$ punktowo.
Zatem, biorąc niezależnie supremum po $g_i$,
\begin{equation}\label{eq:riesz-lower-sum-21}
 \delta\sum_{i=1}^{N}\mu(U_i)\le L(f).
\end{equation}
Natomiast dla $y=f(x)$ równego zero obie sumy
progowe są zerowe. Jeżeli $j\delta<y\le(j+1)\delta$
dla $0\le j<N$, to dokładnie $j$ progów spośród
$\delta,2\delta,\ldots,N\delta$ i dokładnie $j+1$
progów spośród $0,\delta,\ldots,(N-1)\delta$
jest mniejszych od $y$. Otrzymujemy punktowe
oszacowanie przez dolną i górną sumę progową;
po całkowaniu względem $\mu$ daje ono
\[
 \delta\sum_{i=1}^{N}\mu(U_i)
 \le\int_Kf(x)\,\mathrm d\mu(x)
 \le\delta\sum_{i=0}^{N-1}\mu(U_i).
\]
Różnica między lewą i prawą sumą wynosi
$\delta\mu(U_0)\le\delta L(\chi_K)$.
Z \eqref{eq:riesz-upper-sum-21} i
\eqref{eq:riesz-lower-sum-21} wynika
$|L(f)-\int_K f\,\mathrm d\mu|\le2\delta L(\chi_K)$.
Przechodząc do granicy przy $N\to\infty$, otrzymujemy reprezentację
$f\ge0$. Ogólny $f$ rozkładamy na części dodatnią
i ujemną.

\emph{Krok 7: jednoznaczność.}
Dla dowolnej miary regularnej $\rho$ reprezentującej $L$
i otwartego $U$ regularność oraz funkcje odcinające
dają $\rho(U)=a(U)$ (por.
\eqref{eq:cc-open-sup-20}). Stąd z regularności
zewnętrznej $\rho=\mu$ na wszystkich zbiorach
borelowskich.
\end{proof}

\section{Przestrzenie lokalnie zwarte: lokalizacja i sklejanie}
\begin{lemat}[Wyczerpanie przez zbiory relatywnie zwarte]
\label{lem:exhaustion-21}
Niech $X$ będzie lokalnie zwartą, $\sigma$-zwartą przestrzenią Hausdorffa. Istnieją otwarte zbiory $U_n$ takie, że
$U_n\subseteq\overline U_n\subseteq U_{n+1}$,
$\overline U_n$ jest zwarty i $\bigcup_nU_n=X$.
Ponadto istnieją $\chi_n\in C_c(X)$ spełniające
$0\le\chi_n\le1$, $\chi_n=1$ na $\overline U_n$
oraz $\operatorname{supp}\chi_n\subseteq U_{n+1}$.
\end{lemat}
\begin{proof}
Z $\sigma$-zwartości wybieramy zwarte zbiory $D_n$ pokrywające $X$. Kładziemy $U_0=\varnothing$. Jeśli $U_{n-1}$ ma zwarte domknięcie, to $H_n=\overline U_{n-1}\cup D_n$ jest zwarty. Wybieramy skończone pokrycie $H_n$ otwartymi otoczeniami o zwartych domknięciach i za $U_n$ przyjmujemy ich sumę (dla pustego $H_n$ przyjmujemy $U_n=\varnothing$). Domknięcie skończonej sumy jest sumą domknięć, więc $\overline U_n$ jest zwarte. Ponadto $\overline U_{n-1}\subset U_n$ i $D_n\subset U_n$, skąd $\bigcup_n U_n=X$. W szczególności każda lokalnie zwarta ośrodkowa przestrzeń metryczna jest $\sigma$-zwarta: własność Lindelöfa (por. \atlasref{GT-09::thm:topog-metryczne-trzy-warunki}{GT-09: własność Lindelöfa przestrzeni metrycznej ośrodkowej}) daje przeliczalne pokrycie relatywnie zwartymi otoczeniami, a ich domknięcia tworzą pokrycie zwarte.
Funkcje $\chi_n$ istnieją na mocy lematu
\ref{lem:cutoff-radon-20}.
\end{proof}

\begin{proof}[Alternatywny dowód twierdzenia~\ref{thm:riesz-positive-21} w przypadku metrycznym]
Zakładamy w tym dowodzie, że $X$ jest lokalnie zwartą ośrodkową przestrzenią metryczną.
Jeśli $X=\varnothing$, teza zachodzi dla miary zerowej.
\emph{Krok 1: dodatniość lokalna i zwartość.}
Dla dowolnego zwartego $H\subseteq X$ wybieramy
$\eta\in C_c(X)$, $\eta\ge0$, $\eta=1$ na $H$.
Jeśli $\operatorname{supp}f\subseteq H$, to
$-\|f\|_\infty\eta\le f\le\|f\|_\infty\eta$,
więc $|L(f)|\le\|f\|_\infty L(\eta)$.
Wybierzmy wyczerpanie i funkcje $\chi_n$ z lematu
\ref{lem:exhaustion-21}; oznaczmy $K_n=\overline U_n$
i $J_n=\overline U_{n+1}$.
Dla $h\in C(J_n)$ funkcję $\chi_nh$ na $J_n$
przedłużamy przez zero na $X\setminus J_n$.
Jest ona ciągła, ponieważ
$\operatorname{supp}\chi_n\subseteq U_{n+1}$.
Definiujemy $L_n(h)=L(\chi_nh)$; jest to dodatni
funkcjonał na $C(J_n)$.
Jeżeli $J_n=\varnothing$, przyjmujemy miarę zerową.
W przeciwnym razie z twierdzenia~\ref{thm:compact-riesz-21}
istnieje skończona regularna miara $\rho_n$
na $J_n$ reprezentująca $L_n$. Przedłużamy ją
przez zero na dopełnienie $J_n$, otrzymując miarę
borelowską na $X$ (także oznaczaną $\rho_n$).

\emph{Krok 2: zgodność miar na wyczerpaniu.}
Jeżeli $m\ge n$, to $\chi_m=\chi_n=1$ na $U_n$.
Dla dowolnego $f\in C_c(U_n)$ zachodzi zatem
\[
 \int_Xf\,\mathrm d\rho_m=L(f)
  =\int_Xf\,\mathrm d\rho_n.
\]
Przedłużone miary $\rho_j$ są skończonymi miarami
borelowskimi na $X$, więc z twierdzenia
\ref{thm:lch-secondcountable-radon-20} są miarami Radona.
Ich ograniczenia do $U_n$ są skończonymi miarami
borelowskimi na lokalnie zwartej ośrodkowej przestrzeni
metrycznej $U_n$, więc z tego samego twierdzenia są
miarami Radona.
Jednoznaczność z twierdzenia~\ref{thm:uniqueness-cc-20}
daje $\rho_m|_{U_n}=\rho_n|_{U_n}$.


\emph{Krok 3: sklejenie.}
Stawiamy $U_0=\varnothing$, $B_n=U_n\setminus U_{n-1}$.
Dla borelowskiego $E\subseteq X$ określamy
\begin{equation}\label{eq:gluing-measures-21}
 \mu(E)=\sum_{n\ge1}\rho_n(E\cap B_n).
\end{equation}
Ponieważ $B_n$ są parami rozłączne, a suma miar jest
miarą, $\mu$ jest miarą borelowską. Dla $E\subseteq U_m$
przeliczalna addytywność i krok 2 dają
\[
 \mu(E)=\sum_{n=1}^m\rho_n(E\cap B_n)
 =\sum_{n=1}^m\rho_m(E\cap B_n)=\rho_m(E),
\]
ponieważ $B_n\subseteq U_n$ i $E\cap B_n=\varnothing$
dla $n>m$. Dotyczy to wszystkich $E\in\mathcal B(U_m)$,
które są borelowskie również w $X$, gdyż $U_m$ jest otwarty. W szczególności
$\mu(U_m)=\rho_m(U_m)<\infty$.
Zatem $\mu$ jest lokalnie skończona. Rozdział
\ref{ch:radon-measures-20} pokazuje, że na lokalnie
zwartej ośrodkowej przestrzeni metrycznej jest ona
miarą Radona.

\emph{Krok 4: reprezentacja.}
Jeśli $f\in C_c(X)$, jej nośnik mieści się w pewnym
$U_m$ (ponieważ $(U_m)$ jest rosnącym pokryciem
zwartego nośnika). Wtedy
\[
 \int_X f(x)\,\mathrm d\mu(x)
 =\int_X f(x)\,\mathrm d\rho_m(x)
 =L(\chi_m f)=L(f).
\]
Jednoznaczność wynika z twierdzenia
\ref{thm:uniqueness-cc-20}.
\end{proof}

\section{Przestrzeń \texorpdfstring{$C_0(X)$}{C\string_0(X)} oraz miary o skończonej wariacji}
\begin{lemat}[Gęstość $C_c(X)$ w $C_0(X)$]
\label{lem:density-cc-c0-21}
Dla dowolnej lokalnie zwartej przestrzeni Hausdorffa $X$
przestrzeń $C_c(X;\mathbb K)$ jest gęsta w
$C_0(X;\mathbb K)$ w normie supremum, gdzie
$\mathbb K=\mathbb R$ lub $\mathbb C$.
\end{lemat}
\begin{proof}
Dla $f\in C_0(X;\mathbb K)$ i $\varepsilon>0$
zbiór $K=\{|f|\ge\varepsilon\}$ jest zwarty.
Z lematu \ref{lem:cutoff-radon-20} wybieramy
$\eta\in C_c(X)$, $0\le\eta\le1$, $\eta|_K=1$.
Funkcja $\eta f$ ma nośnik zawarty w zwartym
nośniku $\eta$ i jest ciągła. Na $K$ błąd jest
zerowy, poza $K$ mamy $|f-\eta f|\le |f|<\varepsilon$.
Stąd $\|f-\eta f\|_\infty\le\varepsilon$.
\end{proof}

\begin{twierdzenie}[Ograniczone dodatnie funkcjonały na $C_0$]
\label{thm:riesz-c0-positive-21}
Niech $X$ będzie lokalnie zwartą przestrzenią Hausdorffa.
Dla dodatniego liniowego funkcjonału $L$ na $C_c(X;\mathbb R)$
i jego miary reprezentującej $\mu$ następujące warunki
są równoważne: $L$ jest ograniczony w normie supremum;
$\mu(X)<\infty$; $L$ ma ciągłe liniowe rozszerzenie
na $C_0(X;\mathbb R)$. Rozszerzenie jest wtedy
jednoznaczne, dodatnie, dane przez całkę względem $\mu$
i ma normę $\mu(X)$.
Każdy ograniczony dodatni funkcjonał na $C_0(X;\mathbb R)$
powstaje w ten sposób.
\end{twierdzenie}
\begin{proof}
Jeżeli $L$ jest ograniczony, ze wzoru
\eqref{eq:cc-open-sup-20} dla $U=X$ otrzymujemy
$\mu(X)\le\|L\|<\infty$. Nie wymagamy tu
przeliczalnego wyczerpania przestrzeni.
Jeśli $\mu(X)<\infty$, to
$|\int_X f\,\mathrm d\mu|\le\mu(X)\|f\|_\infty$
dla wszystkich $f\in C_0(X;\mathbb R)$.
Całka jest zatem ciągłym dodatnim rozszerzeniem
$L$. Jednoznaczność wynika z gęstości
(lemat \ref{lem:density-cc-c0-21}).
Ograniczenie ciągłego funkcjonału do $C_c(X)$
jest ograniczone, co zamyka cykl implikacji.
Norma rozszerzenia jest nie większa od $\mu(X)$,
a oszacowanie ze wzoru \eqref{eq:cc-open-sup-20}
daje przeciwną nierówność. Dla dowolnego
ograniczonego dodatniego funkcjonału na $C_0(X)$
stosujemy twierdzenie \ref{thm:riesz-positive-21}
do jego ograniczenia do $C_c(X)$, a następnie
korzystamy z jednoznaczności rozszerzenia.
\end{proof}

\begin{lemat}[Aproksymacja względem skończonej miary Radona]
\label{it:lem:radon-cc-l1}
Niech $X$ będzie lokalnie zwartą przestrzenią Hausdorffa.
Jeśli $\sigma$ jest skończoną miarą Radona na $X$, to
$C_c(X;\mathbb K)$ jest gęsta w $L^1(\sigma;\mathbb K)$,
także po uzupełnieniu miary.
Dla mierzalnej funkcji $u$ z $|u|\le1$ można wybrać
przybliżenia $f_n\in C_c(X;\mathbb K)$ spełniające
$\|f_n\|_\infty\le1$.
\end{lemat}
\begin{proof}
Dla miary na $\mathcal B(X)$ gęstość wynika bezpośrednio
z twierdzenia \ref{thm:cc-dense-lp-radon-20} dla $p=1$,
które obowiązuje na dowolnej lokalnie zwartej przestrzeni
Hausdorffa. Niech $\overline\sigma$ oznacza uzupełnienie $\sigma$.
Dla każdego zbioru $E$ mierzalnego względem $\overline\sigma$
istnieje borelowski $B$ taki, że $\overline\sigma(E\triangle B)=0$.
Każdą funkcję prostą względem uzupełnienia można zatem
zastąpić borelowską funkcją prostą równą jej prawie wszędzie,
zastępując w ten sposób wszystkie zbiory występujące
w skończonej reprezentacji. Funkcje proste są gęste w
$L^1(\overline\sigma)$ na mocy twierdzenia \ref{thm:simple-dense-lp-14}.
Aproksymujemy więc najpierw przez taką funkcję prostą,
a następnie przez funkcję z $C_c(X;\mathbb K)$,
stosując nierówność trójkąta dla błędów.
Dla $|u|\le1$ zastępujemy dowolne przybliżenia $g_n$
funkcjami $f_n=P\circ g_n$, gdzie
$P(z)=z/\max(1,|z|)$.
Jest to odwzorowanie ciągłe, $P(0)=0$, więc
$f_n\in C_c(X;\mathbb K)$ i $|f_n|\le1$.
Jeżeli $z=rv$, $r>1$, $|v|=1$, $|w|\le1$, to
\[
 |z-w|^2-|P(z)-w|^2
 =(r-1)(r+1-2\operatorname{Re}(v\overline w))\ge0.
\]
Dla $|z|\le1$ mamy $P(z)=z$.
Stąd $|f_n-u|\le|g_n-u|$, co dowodzi tezy.
\end{proof}

\begin{definicja}[Skończona miara Radona ze znakiem lub zespolona]
\label{idx:21:definicja:08}
Niech $X$ będzie lokalnie zwartą przestrzenią Hausdorffa.
Przeliczalnie addytywną miarę borelowską $\nu$ o wartościach
w $\mathbb K=\mathbb R$ lub $\mathbb C$ nazywamy skończoną
miarą Radona, gdy jej całkowite wahanie $|\nu|$ jest
skończoną miarą Radona. Oznaczamy przestrzeń tych miar
przez $M(X;\mathbb K)$ i kładziemy $\|\nu\|=|\nu|(X)$.
Całkowite wahanie jest określone wzorem
\[
 |\nu|(E)=\sup\left\{\sum_{j=1}^m|\nu(E_j)|:
 E=\bigsqcup_{j=1}^m E_j,\ E_j\in\mathcal B(X)\right\}.
\]
\end{definicja}

\begin{uwaga}[Własności miar ze znakiem i zespolonych]
Korzystamy z wyników teorii miary:
\atlasref{ME-16::thm:jordan-16}{rozkładu Jordana},
\atlasref{ME-16::thm:partition-variation-16}{wzoru na wahanie miary ze znakiem}
oraz \atlasref{ME-16::thm:complex-variation-16}{twierdzenia o wahaniu miary zespolonej}.
Jeżeli $\nu$ jest rzeczywista, to
$|\nu|=\nu^++\nu^-$, gdzie $\nu^+,\nu^-$ są
wzajemnie osobliwymi nieujemnymi częściami Jordana.
Dla miary zespolonej całkę definiujemy liniowo przez
cztery nieujemne części Jordana jej części rzeczywistej
i urojonej. Dla ograniczonej mierzalnej $f$ mamy
\[
 \left|\int_X f\,\mathrm d\nu\right|
 \le\int_X |f|\,\mathrm d|\nu|.
\]
Istotnie, dla funkcji prostej $f=\sum_j a_j\chi_{E_j}$
o rozłącznych zbiorach $E_j$ wynika to z nierówności
$|\nu(E_j)|\le|\nu|(E_j)$ i nierówności trójkąta.
Dla ograniczonych funkcji mierzalnych przechodzimy do
granicy w jednostajnej aproksymacji funkcjami prostymi;
wszystkie cztery miary dodatnie są skończone.
Części Jordana miar $\operatorname{Re}\nu$ i
$\operatorname{Im}\nu$ są zdominowane przez $|\nu|$,
gdyż ich całkowite wahania są nie większe od $|\nu|$
na mocy definicji przez podziały. Są więc Radona
z lematu \ref{it:lem:radon-dominated}. Odwrotnie,
każda kombinacja liniowa skończonych dodatnich miar
Radona jest miarą Radona: jej całkowite wahanie
jest zdominowane przez sumę tych miar pomnożonych
przez moduły współczynników, co wynika z definicji
przez podziały. Lemat \ref{it:lem:radon-dominated}
daje regularność. W szczególności $M(X;\mathbb K)$
jest przestrzenią wektorową. Przyporządkowanie
$\nu\mapsto|\nu|(X)$ jest normą: z definicji przez
podziały wynika $|a\nu|(X)=|a||\nu|(X)$ i
$|\nu+\xi|(X)\le|\nu|(X)+|\xi|(X)$, ponieważ dla
każdego podziału stosujemy nierówność trójkąta
do każdego składnika i bierzemy supremum.
Jeżeli $|\nu|(X)=0$, to dla każdego borelowskiego $E$
mamy $|\nu(E)|\le|\nu|(E)\le|\nu|(X)=0$,
a więc $\nu=0$.
\end{uwaga}

\begin{twierdzenie}[Reprezentacja funkcjonałów rzeczywistych i zespolonych]
\label{thm:riesz-complex-21}
Niech $X$ będzie dowolną lokalnie zwartą przestrzenią
Hausdorffa. Dla $\mathbb K=\mathbb R$ lub $\mathbb C$
każdy ograniczony liniowy funkcjonał $T$ na
$C_0(X;\mathbb K)$ ma dokładnie jedną reprezentację
\[
 T(f)=\int_X f\,\mathrm d\nu,\qquad
 \nu\in M(X;\mathbb K).
\]
Ponadto $\|T\|=|\nu|(X)$, a przyporządkowanie
$\nu\mapsto T_\nu$ jest liniowym izometrycznym
izomorfizmem $M(X;\mathbb K)\cong C_0(X;\mathbb K)'$.
\end{twierdzenie}
\begin{proof}
\emph{Krok 1: rozkład funkcjonału rzeczywistego.}
Dla $f\in C_0(X;\mathbb R)$, $f\ge0$, kładziemy
\[
 T^+(f)=\sup\{T(g):g\in C_0(X;\mathbb R),\ 0\le g\le f\}.
\]
Mamy $0\le T^+(f)\le\|T\|\|f\|_\infty$.
Dodatnia jednorodność wynika z przeskalowania zbioru
dopuszczalnych $g$ (dla skalaru zero wprost).
Dla $f_1,f_2\ge0$ sumy dopuszczalnych $g_1,g_2$
są dopuszczalne dla $f_1+f_2$, więc
$T^+(f_1+f_2)\ge T^+(f_1)+T^+(f_2)$.
Odwrotnie, jeśli $0\le g\le f_1+f_2$, kładziemy
$g_1=g\wedge f_1$, $g_2=g-g_1$. Są to nieujemne
funkcje z $C_0(X)$, $g_1\le f_1$, $g_2\le f_2$.
Własność znikania w nieskończoności wynika z
domkniętości zbiorów poziomicowych i dominacji.
Zatem $T(g)\le T^+(f_1)+T^+(f_2)$.
Biorąc supremum po $g$, otrzymujemy nierówność
przeciwną, a zatem addytywność $T^+$ na stożku.

Definiujemy $T^+(u-v)=T^+(u)-T^+(v)$ dla $u,v\ge0$.
Jeżeli $u-v=u'-v'$, to $u+v'=u'+v$, więc
addytywność dowodzi niezależności od rozkładu.
Dodawanie rozkładów i skalowanie (dla ujemnego
skalaru zamiana składników) dają liniowość.
Otrzymujemy funkcjonał dodatni, a
$|T^+(f)|\le T^+(|f|)\le\|T\|\|f\|_\infty$.
Funkcjonał $T^-=T^+-T$ jest dodatni, bo dla $f\ge0$
można wybrać $g=f$ w supremum. Jest też ograniczony
jako różnica ograniczonych funkcjonałów.
Twierdzenie \ref{thm:riesz-c0-positive-21} daje
skończone miary dodatnie $\alpha,\beta$ reprezentujące
$T^+,T^-$. Stąd $\nu=\alpha-\beta$ reprezentuje $T$
i jest miarą Radona z lematu \ref{it:lem:radon-dominated}.

Sprawdzimy jeszcze, że $\alpha,\beta$ są dokładnie
częściami Jordana $\nu^+,\nu^-$, a nie tylko pewnymi
miarami dodatnimi, których różnicą jest $\nu$.
Niech $P$ będzie dodatnią częścią borelowskiego
rozkładu Hahna miary $\nu$.
Dla $f\in C_0(X;\mathbb R)$, $f\ge0$, oraz
$0\le g\le f$ mamy
\[
 T(g)=\int_X g\,\mathrm d\nu
 \le\int_X f\,\mathrm d\nu^+.
\]
Stąd $T^+(f)\le\int_X f\,\mathrm d\nu^+$.
Z lematu \ref{it:lem:radon-cc-l1} zastosowanego
do skończonej miary Radona $|\nu|$ otrzymujemy
$h_n\in C_c(X;\mathbb R)$ z
$\|h_n-\chi_P\|_{L^1(|\nu|)}\to0$.
Obcinając $h_n$ do przedziału $[0,1]$ przez
$h_n\mapsto0\vee(h_n\wedge1)$, zachowujemy
ciągłość, zwarty nośnik i nie zwiększamy błędu
względem $\chi_P$. Możemy więc założyć
$0\le h_n\le1$. Funkcje $g_n=fh_n$ są
dopuszczalne w definicji $T^+(f)$ i
\[
 \left|\int_X g_n\,\mathrm d\nu
             -\int_X f\chi_P\,\mathrm d\nu\right|
 \le\|f\|_\infty\|h_n-\chi_P\|_{L^1(|\nu|)}
 \longrightarrow0.
\]
Ponieważ $\int_X f\chi_P\,\mathrm d\nu
=\int_X f\,\mathrm d\nu^+$, otrzymujemy równość
$T^+(f)=\int_X f\,\mathrm d\nu^+$.
Części Jordana są miarami Radona przez dominację
przez $|\nu|$. Liniowość i rozkład funkcji na części
dodatnią i ujemną dają równość całek dla wszystkich
funkcji rzeczywistych z $C_0(X)$. Z jednoznaczności
reprezentacji dodatniej
(twierdzenie \ref{thm:riesz-positive-21}) mamy
$\alpha=\nu^+$, a z $\nu=\alpha-\beta$ wynika
$\beta=\nu^-$.

\emph{Krok 2: przypadek zespolony.}
Na rzeczywistych funkcjach z $C_0(X)$ definiujemy
$T_1(f)=\operatorname{Re}T(f)$ i
$T_2(f)=\operatorname{Im}T(f)$. Są rzeczywistymi
ograniczonymi funkcjonałami liniowymi, więc krok 1
daje reprezentujące je miary rzeczywiste $\nu_1,\nu_2$.
Miara $\nu=\nu_1+i\nu_2$ jest Radona dzięki dominacji
$|\nu|\le|\nu_1|+|\nu_2|$. Dla $f=u+iv$ z $u,v$
rzeczywistymi mamy
\[
 T(f)=T_1(u)-T_2(v)+i\bigl(T_2(u)+T_1(v)\bigr)
 =\int_X f\,\mathrm d\nu.
\]

\emph{Krok 3: norma.}
Dla każdej miary $\nu\in M(X;\mathbb K)$ oszacowanie
całki daje $\|T_\nu\|\le|\nu|(X)$.
Dla $\varepsilon>0$ wybieramy skończony borelowski
podział $X=\bigsqcup_{j=1}^m E_j$ z
$\sum_j|\nu(E_j)|>|\nu|(X)-\varepsilon$.
Niech $c_j=\overline{\nu(E_j)}/|\nu(E_j)|$ gdy
$\nu(E_j)\ne0$, a $c_j=0$ w przeciwnym razie.
Funkcja $u=\sum_jc_j\chi_{E_j}$ ma moduł co najwyżej
$1$ i $\int_Xu\,\mathrm d\nu=\sum_j|\nu(E_j)|$.
Dla miary rzeczywistej współczynniki są rzeczywiste.
Lemat \ref{it:lem:radon-cc-l1} dla $\sigma=|\nu|$
daje funkcje $f_n\in C_c(X;\mathbb K)$ o normie
supremum co najwyżej $1$ i
$\int_X|f_n-u|\,\mathrm d|\nu|\to0$. Zatem
\[
 |\nu|(X)-\varepsilon
 <\left|\int_Xu\,\mathrm d\nu\right|
 =\lim_n|T_\nu(f_n)|\le\|T_\nu\|.
\]
Dowolność $\varepsilon$ daje równość norm.
Nie korzystaliśmy z twierdzenia Radona--Nikodyma.

\emph{Krok 4: jednoznaczność i izomorfizm.}
Każda skończona miara Radona definiuje ograniczony
funkcjonał przez całkę. Jeżeli dwie miary definiują
ten sam funkcjonał, ich różnica $\xi$ jest Radona
i $T_\xi=0$. Z kroku 3 mamy $|\xi|(X)=0$,
więc $|\xi(E)|\le|\xi|(X)=0$ dla każdego $E$.
To dowodzi jednoznaczności. Liniowość względem miary wynika najpierw dla
funkcji prostych z równości
$\int_X s\,\mathrm d\nu=\sum_j a_j\nu(E_j)$
dla $s=\sum_j a_j\chi_{E_j}$ na rozłącznych zbiorach,
a następnie dla funkcji ograniczonych przez jednostajną
aproksymację funkcjami prostymi i oszacowanie całki
przez całkowite wahanie każdej z rozważanych miar.
Surjektywność wynika z kroków 1--2, a izometryczność
z kroku 3.
\end{proof}

\begin{proof}[Alternatywny dowód w przypadku zespolonym]
Dla nieujemnego $g\in C_c(X;\mathbb R)$ zdefiniujmy
\[
 M(g)=\sup\{|T(h)|:h\in C_c(X;\mathbb C),\ |h|\le g\}.
\]
Jest to skończone, bo $M(g)\le\|T\|\|g\|_\infty$.
Dodatnia jednorodność $M(ag)=aM(g)$ dla $a>0$ wynika
z bijekcji $h\mapsto ah$ pomiędzy zbiorami dopuszczalnymi;
dla $a=0$ jedyną funkcją dopuszczalną jest zero.
Wykażemy addytywność $M(g_1+g_2)=M(g_1)+M(g_2)$.
Dla nierówności „$\le$” rozkładamy dowolną $h$,
$|h|\le g_1+g_2$, na $h_i=h g_i/(g_1+g_2)$ na
$\{g_1+g_2>0\}$, a w pozostałych punktach przyjmujemy
$h_i=0$. Ciągłość przy zerach mianownika wynika
z $|h_i|\le g_i$; nośnik $h_i$ zawiera się w zwartym
nośniku $g_i$, więc $h_i\in C_c(X;\mathbb C)$.
Mamy $h_1+h_2=h$, stąd
$|T(h)|\le |T(h_1)|+|T(h_2)|\le M(g_1)+M(g_2)$.
Supremum po $h$ daje pierwszą nierówność.
Dla nierówności „$\ge$” ustalmy $\varepsilon>0$ i wybierzmy
$h_i$ z $|h_i|\le g_i$ oraz $|T(h_i)|>M(g_i)-\varepsilon/2$.
Mnożymy $h_i$ przez $\overline{T(h_i)}/|T(h_i)|$,
gdy $T(h_i)\ne0$, a przez $1$ w przeciwnym przypadku.
Nie zmieniamy modułów funkcji, a dzięki liniowości zespolonej
otrzymujemy $T(h_i)=|T(h_i)|$. Wtedy $|h_1+h_2|\le g_1+g_2$
i $M(g_1+g_2)\ge T(h_1+h_2)>M(g_1)+M(g_2)-\varepsilon$.
Dowolność $\varepsilon$ kończy dowód addytywności.
Z addytywności i dodatniej jednorodności $M$ przedłuża
się jednoznacznie do dodatniego rzeczywistego funkcjonału
na $C_c(X;\mathbb R)$: $M(f)=M(f^+)-M(f^-)$.
Dla dwóch rozkładów $f=u-v=u'-v'$ na funkcje nieujemne
mamy $u+v'=u'+v$; addytywność zapewnia niezależność
różnicy $M(u)-M(v)$ od rozkładu. Dodawanie i skalowanie
rozkładów dają liniowość przedłużenia. Z dodatniości
mamy $|M(f)|\le M(|f|)\le\|T\|\|f\|_\infty$,
więc przedłużony funkcjonał jest ograniczony i rozciąga się
na $C_0(X;\mathbb R)$ na mocy lematu
\ref{lem:density-cc-c0-21}.
Twierdzenie~\ref{thm:riesz-c0-positive-21} daje
skończoną miarę Radona $\sigma$ taką, że
$M(g)=\int_Xg\,\mathrm d\sigma$ dla $g\ge0$.
W szczególności dla $h\in C_c(X;\mathbb C)$ zachodzi
\begin{equation}\label{eq:riesz-domination-21}
 |T(h)|\le \int_X|h(x)|\,\mathrm d\sigma(x).
\end{equation}
Z \eqref{eq:riesz-domination-21} wynika, że jeśli $h_1=h_2$
$\sigma$-prawie wszędzie, to $T(h_1-h_2)=0$.
Funkcjonał jest zatem dobrze określony na klasach
równoważności funkcji z $C_c(X;\mathbb C)$ i ograniczony
w normie $L^1(\sigma)$. Klasy te są gęste w
$L^1(\sigma;\mathbb C)$ na mocy lematu
\ref{it:lem:radon-cc-l1}. Zatem $T$ jednoznacznie
przedłuża się do funkcjonału na $L^1(\sigma)$
o normie co najwyżej $1$.
Twierdzenie \ref{thm:dual-l1-15} o dualności
$L^1$--$L^\infty$ (miara $\sigma$ jest skończona) daje $w\in L^\infty(\sigma)$,
$|w|\le1$ $\sigma$-prawie wszędzie, takie że
$T(h)=\int_X h w\,\mathrm d\sigma$.
Kładziemy $\nu(E)=\int_E w\,\mathrm d\sigma$.
Ponieważ $\sigma$ jest skończona i $w\in L^\infty(\sigma)$,
mamy $w\in L^1(\sigma)$. Z twierdzenia \ref{thm:density-variation-16}
wynika, że $\nu$ jest miarą zespoloną oraz
$|\nu|(E)=\int_E|w|\,\mathrm d\sigma\le\sigma(E)$.
Jest to więc skończona miara Radona, gdyż
regularność przechodzi na miary zdominowane (lemat \ref{it:lem:radon-dominated}).
Reprezentacja obowiązuje dla $C_c$, a przez gęstość
także dla $C_0$.

Pozostaje identyczność norm. Nierówność
$\|T\|\le|\nu|(X)$ jest bezpośrednia. Dla odwrotnej
korzystamy z twierdzenia \ref{thm:polar-complex-measure-18}:
jeśli $v$ jest gęstością $\nu$ względem $|\nu|$, to
$|v|=1$ $|\nu|$-prawie wszędzie. Wybieramy mierzalną
funkcję $u=\overline v$ na $\{|v|=1\}$ i $u=0$ poza tym
zbiorem. Mamy $|u|\le1$ oraz $u\,\mathrm d\nu=\mathrm d|\nu|$.
Z regularności $|\nu|$ i aproksymacji funkcji
mierzalnych przez $C_c$ w $L^1(|\nu|)$ istnieją
$f_n\in C_c(X;\mathbb C)$, $\|f_n\|_\infty\le1$,
$f_n\to u$ w $L^1(|\nu|)$. Ograniczenie normy
zapewnia obcięcie aproksymantów do domkniętego
dysku jednostkowego, które nie zwiększa błędu
punktowego względem $u$. Stąd
$T(f_n)\to\int_Xu\,\mathrm d\nu=|\nu|(X)$,
więc $\|T\|\ge|\nu|(X)$.
Dla jednoznaczności niech $\nu,\widetilde\nu$ reprezentują
$T$ i niech $\xi=\nu-\widetilde\nu$. Miara $\xi$ jest
skończoną miarą Radona, a jej funkcjonał całkowy jest zerowy.
Udowodniona równość norm dla funkcjonału całkowego
(także dla $\xi$) daje $|\xi|(X)=0$. Zatem dla każdego
borelowskiego $E$ mamy $|\xi(E)|\le|\xi|(X)=0$,
czyli $\nu=\widetilde\nu$.
\end{proof}



\begin{przyklad}[Funkcjonał ewaluacji]
Dla ustalonego $x_0\in X$ funkcjonał $T(f)=f(x_0)$
na $C_0(X)$ jest ciągły i ma normę $1$.
Jego miarą reprezentującą jest $\delta_{x_0}$.
\end{przyklad}

\begin{przyklad}[Miara nieskończona i funkcjonał lokalny]
Dla $X=\mathbb R$ i $L(f)=\int_{\mathbb R} f(x)
\,\mathrm d\lambda_1(x)$, $f\in C_c(\mathbb R)$,
otrzymujemy dodatni funkcjonał nieograniczony w normie
supremum na całym $C_c$: funkcje odcinające równe $1$
na $[-n,n]$ mają normę supremum $1$, ale całki
co najmniej $2n$. Dlatego $L$ nie przedłuża się
ciągle do $C_0(\mathbb R)$.
\end{przyklad}

\begin{przyklad}[Gęstość i atom w jednym funkcjonale]
Na $C_c(\mathbb R)$ definiujemy $\Lambda(\varphi)=\int_{[0,1]}\varphi(x)\,\mathrm d\lambda_1(x)+2\varphi(0)$. Jest to dodatni funkcjonał liniowy. Miara $\mu(E)=\lambda_1(E\cap[0,1])+2\delta_0(E)$ jest skończoną miarą borelowską, więc z twierdzenia \ref{thm:lch-secondcountable-radon-20} jest Radona. Całka względem sumy miar jest sumą całek, najpierw dla funkcji prostych, potem przez aproksymację monotoniczną i rozkład na części dodatnią i ujemną. Stąd $\mu$ reprezentuje $\Lambda$ i
\[
 \int_{\mathbb R}1\,\mathrm d\mu=3,\quad
 \int_\mathbb R x^2\,\mathrm d\mu(x)=\int_0^1 x^2\,\mathrm dx+2\cdot0^2=1/3,\quad
 \int_\mathbb R\chi_{\{0\}}\,\mathrm d\mu=2.
\]
W rachunku z $x^2$ użyliśmy zgodności całek Riemanna
i Lebesgue'a (twierdzenie \atlasref{ME-08::thm:riemann-to-lebesgue}{ME-08: zgodność całek Riemanna i Lebesgue'a}).
Wszystkie trzy funkcje są całkowalne względem $\mu$,
choć nie należą do $C_c(\mathbb R)$: miara jest skupiona na
$[0,1]$, na którym są ograniczone. Rachunki dotyczą rozszerzonej całki.
\end{przyklad}
\begin{przyklad}[Całka skupiona na okręgu]
Dla $\varphi\in C_c(\mathbb R^2)$ definiujemy $\Lambda(\varphi)=\int_0^{2\pi}\varphi(\cos t,\sin t)\,\mathrm dt$. Funkcjonał jest dodatni i liniowy. Niech $\gamma(t)=(\cos t,\sin t)$ i $\mu=\gamma_*(\lambda_1|_{[0,2\pi]})$. Jest to skończona miara borelowska na $\mathbb R^2$, więc Radona z twierdzenia \ref{thm:lch-secondcountable-radon-20}. Twierdzenie \atlasref{ME-13::thm:pushforward-integral-13}{ME-13: całka względem miary obrazowej}
o całkowaniu względem miary obrazowej oraz twierdzenie
\atlasref{ME-08::thm:riemann-to-lebesgue}{ME-08: zgodność całek Riemanna i Lebesgue'a} o zgodności całki Riemanna
z Lebesgue'a dają reprezentację $\Lambda$. Miara opisuje długość łuku, gdyż $|\gamma'(t)|=1$, a podwójnie przyjmowany punkt końcowy odpowiada zbiorowi parametrów miary zero. Mamy $\mu(S^1)=2\pi$ oraz
\[
 \int_{S^1}x_1^2\,\mathrm d\mu(x)=\int_0^{2\pi}\cos^2t\,\mathrm dt
 =\tfrac12\int_0^{2\pi}(1+\cos(2t))\,\mathrm dt=\pi.
\]
Natomiast $\lambda_2(S^1)=0$. Okrąg jest domknięty
w $\mathbb R^2$, a więc należy do produktowej
$\sigma$-algebry borelowskiej, gdyż $\mathbb R$
ma przeliczalną bazę i otwarte podzbiory produktu
są przeliczalnymi sumami otwartych prostokątów.
Każdy pionowy przekrój okręgu ma co najwyżej dwa punkty,
a więc jednowymiarową miarę Lebesgue'a zero.
Z twierdzeń \atlasref{ME-12::thm:tonelli-12}{ME-12: twierdzenie Tonellego} i
\atlasref{ME-12::thm:lebesgue-products-12}{ME-12: produkty miar Lebesgue'a} otrzymujemy
\[
 \lambda_2(S^1)
 =\int_{\mathbb R}\lambda_1(\{t:(x_1,t)\in S^1\})
      \,\mathrm d\lambda_1(x_1)=0.
\]
\end{przyklad}
\begin{przyklad}[Skończona suma ewaluacji punktowych]
Niech $X$ będzie lokalnie zwartą przestrzenią Hausdorffa. Dla punktów $x_1,\dots,x_m\in X$ oraz liczb $c_1,\dots,c_m>0$ funkcjonał $\Lambda(f)=\sum_{j=1}^m c_jf(x_j)$ na $C_c(X)$ jest dodatni i liniowy. Reprezentuje go miara Radona $\mu=\sum_{j=1}^m c_j\delta_{x_j}$. Istotnie, dla dowolnego borelowskiego $E$ zwarty zbiór $\{x_j:x_j\in E\}$ jest podzbiorem $E$ o tej samej mierze, a otwarty zbiór $X\setminus\{x_j:x_j\notin E\}$ jest nadzbiorem $E$ o tej samej mierze. To daje obie regularności; skończoność wynika z sumy wag. Reprezentacja wynika z definicji całki względem mas Diraca i addytywności względem miary. Twierdzenie \ref{tw:radon-unikalnosc} pokazuje, że żadna inna regularna miara borelowska nie daje tych samych wartości $\Lambda(f)$ dla wszystkich $f\in C_c(X)$.
\end{przyklad}

\begin{wniosek}[Porządek dodatnich funkcjonałów i miar]
\label{cor:riesz-order-21}
Niech $X$ będzie lokalnie zwartą przestrzenią
Hausdorffa, a $L_1,L_2:C_c(X;\mathbb R)\to\mathbb R$
dodatnimi funkcjonałami liniowymi, reprezentowanymi
odpowiednio przez miary Radona $\mu_1,\mu_2$.
Wówczas następujące warunki są równoważne:
\begin{enumerate}[label=(\roman*)]
\item $L_1(f)\le L_2(f)$ dla każdego nieujemnego $f\in C_c(X)$;
\item $\mu_1(E)\le\mu_2(E)$ dla każdego zbioru borelowskiego
$E\subseteq X$.
\end{enumerate}
\end{wniosek}
\begin{proof}
Implikacja (ii)$\Rightarrow$(i) wynika z monotoniczności całki.
Załóżmy (i). Dla otwartego $U\subseteq X$ zachodzi
\[
 \mu_1(U)=\sup\{L_1(f):f\in C_c(X),\ 0\le f\le1,
 \ \operatorname{supp}f\subseteq U\}.
\]
Wzór ten udowodniono w rozdziale~\ref{ch:radon-measures-20};
obejmuje także przypadek $\mu_1(U)=+\infty$.
Nierówność dla funkcjonałów daje
$\mu_1(U)\le\mu_2(U)$. Dla dowolnego borelowskiego
$E$ z regularności zewnętrznej otrzymujemy
$\mu_1(E)=\inf_{U\supseteq E}\mu_1(U)
\le\inf_{U\supseteq E}\mu_2(U)=\mu_2(E)$.
\end{proof}

\begin{uwaga}[Gdzie użyliśmy twierdzenia Tichonowa?]
W konstrukcji na zwartej przestrzeni metrycznej nie
potrzebowaliśmy twierdzenia Tichonowa: miarę zewnętrzną
zbudowaliśmy bezpośrednio z dodatniego funkcjonału.
W rozdziale~\ref{ch:haar-23}, poświęconym mierze Haara,
pojawi się natomiast zwartość produktu nieprzeliczalnie
wielu przedziałów.
Nie należy jej mylić ze zwartą metryczną przestrzenią,
dla której wystarczy argument diagonalny.
\end{uwaga}

\section{Konsekwencje: oddzielenie miar i problem momentów}
\begin{wniosek}[Miary odróżniane przez funkcje ciągłe]
\label{cor:measures-separate-21}
Niech $\mu,\nu$ będą miarami Radona na lokalnie zwartej
przestrzeni Hausdorffa $X$. Wtedy
$\mu=\nu$ wtedy i tylko wtedy, gdy
$\int_Xf\,\mathrm d\mu=\int_Xf\,\mathrm d\nu$
dla wszystkich $f\in C_c(X;\mathbb R)$.
\end{wniosek}
\begin{proof}
Implikacja w prawo jest oczywista, a odwrotna to
\ref{thm:uniqueness-cc-20}.
\end{proof}

\begin{uwaga}[Problem momentów jako ostrzeżenie]
Jeżeli na $\mathbb R$ znamy całki
$\int_{\mathbb R}x^k\,\mathrm d\mu(x)$ dla wszystkich
$k\ge0$, nie zawsze wyznacza to jednoznacznie miarę.
Funkcje wielomianowe nie należą do $C_0(\mathbb R)$
i nie ma automatycznej gęstości wielomianów w normie
supremum na całej prostej. Na zwartym przedziale
sytuacja jest inna: twierdzenie Stone'a--Weierstrassa
przybliża dowolną funkcję ciągłą wielomianami,
więc momenty jednoznacznie wyznaczają skończoną miarę.
\end{uwaga}

\begin{przyklad}[Norma funkcjonału atomowego]
\label{ex:riesz-complex-atoms-21}
Niech $X$ będzie lokalnie zwartą przestrzenią Hausdorffa,
$x_1,\ldots,x_m\in X$ parami różnymi punktami,
a $a_1,\ldots,a_m\in\mathbb C$. Na $C_0(X;\mathbb C)$
określmy funkcjonał
\[
 T(f)=\sum_{j=1}^ma_jf(x_j).
\]
Reprezentuje go miara zespolona
$\nu=\sum_{j=1}^ma_j\delta_{x_j}$, a jej wariacja
całkowita wynosi
$|\nu|=\sum_{j=1}^m|a_j|\delta_{x_j}$.
W szczególności
$\|T\|=\sum_{j=1}^m|a_j|$.

Sprawdźmy bezpośrednio ostatnią równość. Z nierówności
trójkąta otrzymujemy
$|T(f)|\le\|f\|_\infty\sum_j|a_j|$.
Ponieważ punktów jest skończenie wiele i są różne,
dla każdej pary punktów wybieramy rozłączne otwarte
otoczenia z aksjomatu Hausdorffa. Dla każdego punktu
przecinamy skończenie wiele przypisanych mu otoczeń.
Otrzymujemy parami rozłączne otwarte otoczenia, które
możemy zmniejszyć do otwartych $U_j$ o zwartych domknięciach
na mocy lematu \atlasref{GT-14::lem:topog-baire-zwarte-domkniecie}{GT-14: otoczenie o zwartym domknięciu}. Z lematu
\ref{lem:cutoff-radon-20} istnieją
$\varphi_j\in C_c(X)$, $0\le\varphi_j\le1$,
$\varphi_j(x_j)=1$ i
$\operatorname{supp}\varphi_j\subseteq U_j$.
Dla $a_j\ne0$ kładziemy
$b_j=\overline{a_j}/|a_j|$, a dla $a_j=0$
przyjmujemy $b_j=0$. Wtedy
$f=\sum_{j=1}^m b_j\varphi_j$ należy do
$C_c(X;\mathbb C)$ i, dzięki rozłączności nośników,
$\|f\|_\infty\le1$. Ponadto
$T(f)=\sum_j a_jb_j=\sum_j|a_j|$,
co dowodzi wzoru na normę.
\end{przyklad}

\section{Podsumowanie i przygotowanie do zbieżności miar}
Reprezentacja dodatnich funkcjonałów pozwala przejść
od analizy na $C_c(X)$ do teorii miar Radona. Dla
funkcjonałów na $C_0(X)$ skończoność normy odpowiada
skończoności masy, zaś dla zespolonych funkcjonałów
norma jest wariacją całkowitą. W rozdziale~\ref{ch:weak-measures-22}
zastosujemy zwarty przypadek twierdzenia do wykazania
ciągowej zwartości przestrzeni prawdopodobieństw na
zwartej przestrzeni metrycznej.
